\section{Iitaka fibrations}


\subsection{Iitaka dimension}

  % \Yang{}

  \begin{definition}
    Let $D$ be a divisor on a projective variety $X$.
    We define the \emph{Iitaka dimension} of $D$ as
    \[ \kappa(X, D) \coloneqq \max \left\{ k \in \bbZ_{\geq 0} \;\middle|\; \limsup_{m \to +\infty} \frac{h^0(X, mD)}{m^k} > 0 \right\} \in \{-\infty\} \cup \bbZ_{\geq 0}. \]
    If $D = K_X$ is the canonical divisor, we call $\kappa(X, K_X)$ the \emph{Kodaira dimension} of $X$ and denote it by $\kappa(X)$ simply.
  \end{definition}
  

  \begin{proposition}
    Let $D$ be an effective divisor on a variety $X$ with fibration $f: X \setminus \bfB(D) \surjmap Y$.
    Then we have $\kappa(X, D) = \dim Y$.
    % \Yang{}
    % Then there exists a fibration $f: X \to Y$ to a quasi-projective variety $Y$ such that $\calO_X(mD) \cong f^* \calO_Y(1)$ for some ample divisor on $Y$, and we have $\kappa(X, D) = \dim Y$.
    % \Yang{}
  \end{proposition}
  \begin{proof}
    \Yang{}
  \end{proof}

  \begin{remark}
    Iitaka dimension is not a numerical invariant, i.e. it may happen that $D \equiv D'$ but $\kappa(X, D) \neq \kappa(X, D')$.
    \Yang{arXiv:1904.10832}
  \end{remark}



\subsection{Iitaka fibrations}

  In characteristic zero, we have the following birational analog of the above proposition, which is the main result of this section.

  \begin{theorem}
    Let $\kk$ be a field of characteristic zero and $X$ a projective variety over $\kk$.
    Let $D$ be a divisor on $X$ with $\kappa(X, D) \geq 0$.
    There exists a fibration $f: X' \surjmap Y$ with $X'$ birational to $X$ such that for sufficiently large and \Yang{} divisible $m > 0$, the fibration $\phi_m: X \dashrightarrow Y_m$ associated to $mD$ is birational to $f$.

    We call such a rational map an \emph{Iitaka fibration} of $D$.
    \Yang{}
  \end{theorem}

    \Yang{}


